\chapter*{PREFÁCIO}
\addcontentsline{toc}{chapter}{PREFÁCIO}

A decisão de escrever esta obra surgiu da constatação de que a Estatística ocupa posição cada vez mais central na estrutura decisória do Estado contemporâneo. Em contextos marcados por restrições orçamentárias, demandas sociais crescentes e exigência de transparência institucional, governar passou a significar decidir sob condições explícitas de incerteza.  \vskip0.3cm

Nesse cenário, a Estatística não pode ser compreendida apenas como conjunto de técnicas matemáticas aplicadas à mensuração de fenômenos. Ela constitui linguagem formal da evidência, instrumento de avaliação e, sobretudo, fundamento da racionalidade administrativa. A qualidade da informação pública influencia diretamente a legitimidade das decisões governamentais. \vskip0.3cm

A experiência acumulada ao longo de duas décadas de atuação na gestão pública estadual revelou que a ausência de cultura analítica consistente compromete diagnósticos, distorce prioridades e fragiliza políticas públicas. Por outro lado, quando corretamente compreendidos, os princípios estatísticos ampliam a capacidade institucional de planejamento, monitoramento e avaliação. \vskip0.3cm

Este livro foi concebido com o propósito de articular fundamentos estatísticos clássicos a aplicações concretas no campo da administração pública. O objetivo não é apenas apresentar conceitos, mas demonstrar como a incorporação sistemática da análise quantitativa contribui para decisões mais responsáveis, transparentes e fundamentadas. \vskip0.3cm

A Estatística fornece instrumentos para definir metas, avaliar desempenho, estimar impactos e interpretar variabilidade. Contudo, também impõe disciplina intelectual: exige explicitação de hipóteses, reconhecimento de margens de erro e consciência dos limites da modelagem. Essa postura crítica é essencial para evitar interpretações simplificadoras e uso inadequado de indicadores. \vskip0.3cm

Embora estruturado como obra introdutória, o texto busca preservar rigor conceitual e estimular reflexão sobre o papel institucional da análise estatística. A formação quantitativa do gestor público não constitui luxo acadêmico, mas requisito de maturidade administrativa. \vskip0.3cm

\newpage
Os exemplos e aplicações apresentados ao longo do livro foram desenvolvidos a partir de experiências técnicas e acadêmicas acumuladas em diferentes contextos da gestão pública. A utilização de ferramentas computacionais de código aberto reforça o compromisso com acessibilidade, replicabilidade e transparência metodológica. \vskip0.3cm

Espera-se que esta obra contribua para fortalecer a cultura de dados no setor público e estimular uma compreensão mais ampla da Estatística como instrumento de governança. Se ao final da leitura o leitor reconhecer que decidir implica administrar incertezas com responsabilidade técnica e ética, então o propósito deste trabalho estará cumprido.

\vskip1cm
\begin{centering}

\textbf{MÁRIO DIEGO ROCHA VALENTE} \\
(Estatístico, Analista de Trânsito/DETRAN-PA) \\

\vskip0.7cm
\textbf{GIOVANA RAIOL PIRES} \\
(Estatístico, Técnico em Gestão Pública/SEPLAD-PA) \\

\vskip0.7cm
\textbf{HÉLITON RIBEIRO TAVARES} \\
(Estatístico, Prof.\ Dr.\ Titular/UFPA) \\

\end{centering}

%\vskip1cm

%PRÉ-LANÇAMENTO DO LIVRO \vskip0.3cm
%Ocorrerá durante a Semana Nacional do Administrador - 2024\\
%(\textbf{Universidade Federal Rural da Amazônia)} \\
%(\textbf{Campus Parauapebas)} \\

%Belém, 23 de fevereiro de 2026 \\


